\documentclass[11pt,a4paper]{article}
\usepackage[utf8]{inputenc}
\usepackage[T1]{fontenc}
\usepackage{lmodern}
\usepackage{textcomp}
\usepackage{amsmath,amssymb}
\usepackage{graphicx}
\usepackage{booktabs}
\usepackage{longtable}
\usepackage{array}
\usepackage{listings}
\usepackage{xcolor}
\usepackage[margin=1in]{geometry}
\usepackage{float}
\usepackage[hidelinks,breaklinks]{hyperref}
\usepackage{microtype}
\setlength{\parskip}{4pt}
\lstset{basicstyle=\ttfamily\scriptsize,breaklines=true,breakatwhitespace=false,
  postbreak=\mbox{\textcolor{gray}{$\hookrightarrow$}\space},frame=single,
  framerule=0.3pt,rulecolor=\color{gray!60},columns=fullflexible,keepspaces=true,
  showstringspaces=false,numbers=none,xleftmargin=4pt,xrightmargin=4pt,
  literate={\_}{{\_}}1}
\graphicspath{{figures/}}
\providecommand{\tightlist}{}
\title{Transcript: Lyman-$\alpha$ Warm Dark Matter Limits Compared with Non-Resonant Sterile-Neutrino Production}
\author{HEP-Genesis-Agent session record \\ \small Project: ramp-combo-lyman-wdm}
\date{2026-10-19}
\begin{document}
\maketitle
\noindent\textit{\small The conversation verbatim, generated mechanically from the session record.}
\vspace{1em}



\appendix
\section{Conversation transcript}
\label{app:transcript}
The exchange verbatim (assistant text as shown to the user; tool calls summarized by name --- see the tool-call log for arguments and results).
\subsection*{User \small(2026-10-19 14:31:14)}
Survey recent Lyman-alpha forest constraints on warm dark matter. How do they compare with sterile-neutrino non-resonant production? Use the sterile-nu tools for production fractions.

\subsection*{Assistant \small(research mode, 2026-10-19 14:31:14)}
\textit{Plan:}\begin{enumerate}
  \item (closed) Use the sterile-nu server to compute non-resonant (Dodelson-Widrow) production fractions Omega\_s/Omega\_dm and related observables across keV masses and mixing angles. Then locate the DW relic-abundance curve (Omega\_s h\textasciicircum{}2 = 0.12) in the (m\_s, sin\textasciicircum{}2 2theta) plane.
  \item (closed) Verify the recent Lyman-alpha thermal-relic WDM limits and their systematic variations from the papers' text, and verify the thermal-relic to DW sterile-neutrino mass conversion.
  \item (closed) Build the X-ray decay and mixed cold-plus-warm side of the comparison. This means X-ray mixing-angle limits at the DW-equivalent masses, the predicted line energies and fluxes, and the published Lyman-alpha bounds on the warm fraction f versus mass.
  \item (blocked) Extract numerical NuSTAR X-ray upper limits on sin\textasciicircum{}2 2theta at m\_s = 20, 29, 41 and 46 keV. Combine them with the DW relic curve to get the X-ray-allowed fraction f\_X(m\_s), compare f\_X with the Boyarsky+2009 mixed cold-plus-warm bound, and give a final verdict on pure non-resonant production over 20-46 keV.
  \item (closed) Cross-check the sterile-nu server's DW yield normalization and its sin\textasciicircum{}2 2theta scaling against published production calculations. Derive a correction factor and propagate it to the relic curve (f-003), the fractions f (f-004) and the Delta N\_eff statement.
  \item (closed) Test how robust the thermal-relic to NRP mass mapping is by tabulating paper-quoted NRP-equivalent masses next to the Viel+2005-derived values. Run one targeted search for post-2023 Lyman-alpha WDM limits from DESI P1D or emulator re-analyses.
\end{enumerate}
\textit{Tool calls:} \texttt{list\_sterile\_nu\_skills}, \texttt{list\_observables}, \texttt{list\_agent\_files}, \texttt{list\_input\_files}, \texttt{list\_datasets}, \texttt{list\_files}, \texttt{get\_dispatch}, \texttt{search\_arxiv\_abstracts}, \texttt{search\_arxiv\_abstracts}, \texttt{search\_arxiv\_abstracts}, \texttt{search\_arxiv}, \texttt{search\_arxiv\_abstracts}, \texttt{search\_arxiv\_abstracts}, \texttt{get\_arxiv\_metadata}, \texttt{get\_arxiv\_metadata}, \texttt{get\_arxiv\_metadata}, \texttt{get\_arxiv\_metadata}, \texttt{get\_arxiv\_metadata}, \texttt{get\_arxiv\_metadata}, \texttt{describe\_sterile\_nu\_tools}, \texttt{list\_sterile\_nu\_skills}, \texttt{load\_sterile\_nu\_skill}, \texttt{predict\_sterile\_nu}, \texttt{download\_arxiv\_full\_text}, \texttt{predict\_sterile\_nu\_grid}, \texttt{download\_arxiv\_full\_text}, \texttt{describe\_sterile\_nu\_tools}, \texttt{predict\_sterile\_nu\_grid}, \texttt{download\_arxiv\_full\_text}, \texttt{load\_table\_data}, \texttt{read\_arxiv\_chunk}, \texttt{read\_arxiv\_chunk}, \texttt{read\_arxiv\_chunk}, \texttt{read\_text\_file}, \texttt{predict\_sterile\_nu\_grid}, \texttt{predict\_sterile\_nu\_xray\_line}, \texttt{predict\_sterile\_nu\_xray\_line}, \texttt{predict\_sterile\_nu\_xray\_line}, \texttt{predict\_sterile\_nu\_xray\_line}, \texttt{predict\_sterile\_nu\_xray\_line}, \texttt{cite}, \texttt{cite}, \texttt{cite}, \texttt{get\_arxiv\_metadata}, \texttt{get\_arxiv\_metadata}, \texttt{get\_arxiv\_metadata}, \texttt{predict\_sterile\_nu\_grid}, \texttt{scan\_sterile\_nu\_parameters}, \texttt{post\_finding}, \texttt{post\_finding}, \texttt{search\_arxiv\_abstracts}, \texttt{search\_arxiv\_abstracts}, \texttt{get\_arxiv\_metadata}, \texttt{get\_arxiv\_metadata}, \texttt{cite}, \texttt{post\_finding}, \texttt{post\_finding}, \texttt{post\_finding}, \texttt{post\_finding}, \texttt{post\_question}, \texttt{plot\_sterile\_nu\_exclusion}, \texttt{search\_arxiv\_abstracts}, \texttt{post\_finding}, \texttt{post\_finding}, \texttt{post\_question}, \texttt{post\_dead\_end}, \texttt{post\_finding}, \texttt{search\_arxiv\_abstracts}, \texttt{download\_arxiv\_full\_text}, \texttt{predict\_sterile\_nu\_xray\_line}, \texttt{predict\_sterile\_nu\_xray\_line}, \texttt{download\_arxiv\_full\_text}, \texttt{get\_arxiv\_metadata}, \texttt{get\_arxiv\_metadata}, \texttt{get\_arxiv\_metadata}, \texttt{read\_arxiv\_chunk}, \texttt{get\_arxiv\_metadata}, \texttt{read\_arxiv\_chunk}, \texttt{read\_arxiv\_chunk}, \texttt{search\_arxiv\_abstracts}, \texttt{read\_arxiv\_chunk}, \texttt{get\_paper\_figures}, \texttt{search\_arxiv\_abstracts}, \texttt{get\_arxiv\_metadata}, \texttt{post\_finding}, \texttt{post\_finding}, \texttt{post\_dead\_end}, \texttt{post\_finding}, \texttt{post\_finding}, \texttt{post\_finding}, \texttt{post\_dead\_end}, \texttt{post\_question}, \texttt{post\_finding}, \texttt{post\_finding}, \texttt{post\_finding}, \texttt{post\_finding}, \texttt{post\_question}\par
\textbf{Survey} --- 52 tools on 5 servers, 7 read-only probes, 12 papers found.

Warning: the 'sterile-nu tools' named in the task (for production fractions) are not among the available tools. Only arXiv and web search tools were provided, so no production fractions or transfer functions were computed here. Those numbers must come from a dedicated Boltzmann or production code, or be taken from the literature. Lyman-alpha forest flux-power-spectrum analyses set lower limits on the thermal-relic WDM mass that rose from about 0.75 keV (2000) to 3.3 keV (Viel+2013, HIRES/MIKE). Later limits are 5.3 keV from XQ-100 plus HIRES/MIKE (Irsic+2017), 5.3 keV from SDSS DR14 (z\textless{}4.5) plus XQ-100 (Palanque-Delabrouille+2019), 3.1 keV from the Cholla suite of 1080 simulations on Boera+2019 data (Villasenor+2022, which weakly prefers about 4.5 keV), and 5.7 keV from high-z HIRES/UVES data with k\_max=0.2 s/km (Irsic+2023). All are 95\% CL. The main things that loosen these limits are the IGM thermal history (a sudden temperature jump weakens Irsic+2017 to 3.5 keV), peculiar-velocity smearing, patchy reionization, and the choice of k\_max (restricting Irsic+2023 to k\textless{}0.1 s/km gives 4.1 keV). For non-resonant (Dodelson-Widrow) sterile neutrinos, the standard conversion from Viel+2005 (astro-ph/0501562) is m\_s $\approx$ 4.46 keV (m\_x/keV)\textasciicircum{}\{4/3\} (omega\_dm/0.1225)\textasciicircum{}\{-1/3\}. This is my own estimate from memory, not checked against the paper text in this session. Under it, the 3.1-5.7 keV relic limits correspond to roughly m\_s ? 20-46 keV. That is far above the mass at which DW sterile neutrinos are excluded by X-ray decay limits (this needs confirming; see 1512.05369 and 1702.08430). So pure DW production as all of the dark matter is excluded, and viable sterile-neutrino dark matter needs resonant production, non-standard histories, or a fraction below 1.

Quantities worth checking: DW-equivalent sterile-neutrino mass for each relic limit (3.1, 3.3, 4.1, 5.3, 5.7 keV), using m\_s $\approx$ 4.46 keV (m\_x/keV)\textasciicircum{}\{4/3\} (omega\_dm/0.1225)\textasciicircum{}\{-1/3\}; this gives roughly 20, 22, 29, 41 and 46 keV and must be checked against the original paper.; Mixing angle sin\textasciicircum{}2(2theta) needed for DW production to give Omega\_s h\textasciicircum{}2 = 0.12 as a function of m\_s, compared with X-ray decay limits.; Fraction f = Omega\_s/Omega\_dm produced non-resonantly at the X-ray-allowed mixing, and the mixed cold-plus-warm Lyman-alpha bound on f at that mass.; Average momentum \textless{}p/T\textgreater{} of the DW spectrum relative to a thermal one (about 0.9 times the Fermi-Dirac value), which sets the free-streaming length.; Half-mode wavenumber k\_1/2 and the suppression at k = 10 h/Mpc for DW versus thermal-relic transfer functions.; How much the relic limit shifts with k\_max and with thermal-history priors (5.7 vs 4.1 keV; 5.3 vs 3.5 keV)..

\textbf{Plan} --- 3 lines of investigation:

\begin{enumerate}
  \item Use the sterile-nu server to compute non-resonant (Dodelson-Widrow) production fractions Omega\_s/Omega\_dm and related observables across keV masses and mixing angles. Then locate the DW relic-abundance curve (Omega\_s h\textasciicircum{}2 = 0.12) in the (m\_s, sin\textasciicircum{}2 2theta) plane. --- The task explicitly asks for production fractions from the sterile-nu tools. The survey wrongly said these tools are absent: they are in the catalog (serve-sbbn.sh). This direction gives the central quantitative comparison against the Lyman-alpha mass limits.
  \item Verify the recent Lyman-alpha thermal-relic WDM limits and their systematic variations from the papers' text, and verify the thermal-relic to DW sterile-neutrino mass conversion. --- The survey numbers (3.1, 3.3, 5.3, 5.7, 4.1 and 3.5 keV) come from abstracts. The conversion m\_s $\approx$ 4.46 keV (m\_x/keV)\textasciicircum{}\{4/3\} (omega\_dm/0.1225)\textasciicircum{}\{-1/3\} is from memory and must be checked before any DW comparison.
  \item Build the X-ray decay and mixed cold-plus-warm side of the comparison. This means X-ray mixing-angle limits at the DW-equivalent masses, the predicted line energies and fluxes, and the published Lyman-alpha bounds on the warm fraction f versus mass. --- The conclusion that pure DW is excluded requires the X-ray upper limit on mixing to lie below the DW relic curve at m\_s ? 20 keV. Fractional DW dark matter then needs the mixed-model Lyman-alpha bound on f, to compare with the fractions from d-01.
Critic's view: revise --- The sterile-nu server is serve-sbbn.sh. Its probes (the eV-scale SBL anomaly check, BBN Yp and D/H, N\_eff) suggest it may only model thermalized light steriles, not keV Dodelson-Widrow freeze-in. Make step 1 a hard gate. If describe\_sterile\_nu\_tools or a single predict\_sterile\_nu call at 7.1 keV does not return a non-resonant relic fraction, stop the grid and scan. Report that the production fractions must come from the literature, and do not interpolate a relic curve from outputs that are out of validity.; Add a robustness cross-check. Compare the tool's DW relic curve at 2-3 masses against a published analytic DW abundance, for example the Abazajian/Asaka fits in 1705.01837. Without this, f values tied only to tool calls cannot be verified. Also note that predict\_sterile\_nu\_grid and scan\_sterile\_nu\_parameters overlap; drop the scan if the grid already gives the plane, and save the calls.; This does not fit 13 calls. Six papers, each needing download, chunk reads and cite, plus a post-2023 search, needs about 20 or more calls. Prioritize Viel+2005 for the conversion and 2309.04533 and 1702.01764 for the limits and systematics, all with full text. Take 2209.14220, 1911.09073 and 1306.2314 from abstracts or metadata, flagged as abstract-level. Drop the optional figure step..

\end{enumerate}
\textbf{Wave 1} --- now exploring, in parallel: Use the sterile-nu server to compute non-resonant (Dodelson-Widrow) production fractions Omega\_s/Omega\_dm and related observables across keV masses and mixing angles. Then locate the DW relic-abundance curve (Omega\_s h\textasciicircum{}2 = 0.12) in the (m\_s, sin\textasciicircum{}2 2theta) plane.; Verify the recent Lyman-alpha thermal-relic WDM limits and their systematic variations from the papers' text, and verify the thermal-relic to DW sterile-neutrino mass conversion.; Build the X-ray decay and mixed cold-plus-warm side of the comparison. This means X-ray mixing-angle limits at the DW-equivalent masses, the predicted line energies and fluxes, and the published Lyman-alpha bounds on the warm fraction f versus mass..

\textbf{Finding:} On the sterile-nu server's DW relic curve (Omega\_s h\textasciicircum{}2 = 0.12), the mixing is sin\textasciicircum{}2 2theta = 2.545e-10, 3.946e-11, 2.022e-11, 1.084e-11 and 8.81e-12 at m\_s = 7.1, 20, 29, 41 and 46 keV. The corresponding radiative lifetimes are 1.6e27, 5.82e25, 1.77e25, 5.85e24 and 4.05e24 s.

\textbf{Finding:} At the DW relic mixing, the line energies are E = 3.55, 10, 14.5, 20.5 and 23 keV for m\_s = 7.1, 20, 29, 41 and 46 keV. The whole-halo M31 fluxes are 1.75e-3, 1.71e-2, 3.88e-2, 8.31e-2 and 0.107 ph/s/cm\textasciicircum{}2. The tool's tabulated limits exclude m\_s = 7.1 keV (NuSTAR M31 Roach+2020 and XMM blank sky Dessert+2020) and m\_s = 20 keV (NuSTAR M31). They report no tabulated limit within 30\% in mass at 29, 41 or 46 keV. That means no coverage, not that these masses are allowed.

\textbf{Finding:} On the sterile-nu server (DW mechanism), Omega\_s h\textasciicircum{}2 scales linearly with sin\textasciicircum{}2 2theta and roughly as m\_s\textasciicircum{}1.8. Setting Omega\_s h\textasciicircum{}2 = 0.12 gives a DW relic curve of sin\textasciicircum{}2 2theta $\approx$ 8.7e-9 (1 keV), 2.5e-9 (2 keV), 1.2e-9 (3 keV), 4.8e-10 (5 keV), 2.5e-10 (7.1 keV), 1.4e-10 (10 keV), 6.6e-11 (15 keV), 3.9e-11 (20 keV), 3.3e-11 (22 keV), 2.0e-11 (29 keV), 1.08e-11 (41 keV), 8.8e-12 (46 keV) and 7.6e-12 (50 keV).

\textbf{Finding:} DW fractions f = Omega\_s h\textasciicircum{}2 / 0.12 from the server: at 7.1 keV and sin\textasciicircum{}2 2theta = 5e-11, f = 0.196 (Omega\_s h\textasciicircum{}2 = 0.0236). At sin\textasciicircum{}2 2theta = 1e-12, f = 0.025 (20 keV), 0.030 (22 keV), 0.049 (29 keV), 0.092 (41 keV) and 0.114 (46 keV). At 1e-11, f is ten times these values; 46 and 50 keV then overclose (Omega\_s h\textasciicircum{}2 = 0.136 and 0.158). The free-streaming length is 151 kpc at 7.1 keV, 53.7 kpc at 20 keV and 107 kpc at 10 keV. Delta N\_eff is at most 3e-4 and BBN is unconstrained for any DW keV dark matter near the relic curve.

\textbf{Finding:} Verified from the text of Viel+2005 (astro-ph/0501562, Eq. 5): m\_s(DW) = 4.43 keV (m\_x/1 keV)\textasciicircum{}\{4/3\} (0.25*0.7\textasciicircum{}2/omega\_x)\textasciicircum{}\{1/3\}, normalized at omega\_x = 0.1225; the survey's 4.46 keV prefactor is not in the paper but equals 4.43*(0.1225/0.12)\textasciicircum{}\{1/3\} = 4.461 keV once rescaled to omega\_dm = 0.12. DW-equivalent masses (omega\_dm=0.1225 / 0.12): m\_x=3.1 -\textgreater{} 20.0/20.2 keV; 3.3 -\textgreater{} 21.8/21.9; 3.5 -\textgreater{} 23.5/23.7; 4.1 -\textgreater{} 29.1/29.3; 4.5 (Villasenor peak) -\textgreater{} 32.9/33.1; 5.3 -\textgreater{} 40.9/41.2; 5.7 -\textgreater{} 45.1/45.4 keV (m\_x\textasciicircum{}\{4/3\} = 4.520, 4.913, 5.314, 6.562, 7.430, 9.241, 10.182). The survey's range of 20-46 keV is confirmed to within 0.6 keV.

\textbf{Finding:} Verified Lyman-alpha thermal-relic WDM limits: Viel+2013 (HIRES/MIKE, z\textgreater{}4) m\_x\textgreater{}3.3 keV (2 sigma; 1, 2, 2.5 keV disfavoured at about 9, 4, 3 sigma); Irsic+2017 XQ-100 alone 1.4 keV, HIRES/MIKE 4.1 keV, combined 5.3 keV (2 sigma), falling to 3.5 keV if sudden temperature jumps of up to 5000 K are allowed; Palanque-Delabrouille+2019 SDSS DR14 z\textless{}4.5 plus XQ-100 m\_X\textgreater{}5.3 keV (95\% CL); Villasenor+2022 (1080 Cholla sims, Boera+2019 data, z=4.0-5.2) m\_WDM\textgreater{}3.1 keV (95\% CL) with a likelihood peak at 4.5 keV and FPS suppression saturating at k\textgreater{}\textasciitilde{}0.1 s/km from peculiar-velocity smearing; Irsic+2023 HIRES/UVES z=4.2-5.0 m\_x\textgreater{}5.7 keV (95\% CL, k\_max=0.2 s/km), falling to 4.1 keV for k\_max\textless{}0.1 s/km with earlier thermal priors, with peculiar velocities and inhomogeneous reionization named as the main small-scale systematics.

\textbf{Finding:} The mixed cold-plus-warm Lyman-alpha bound of Boyarsky+2009 (arXiv:0812.0010), with masses expressed as non-resonantly produced sterile-neutrino masses, allows any mass in 5 keV \textless{} m\_NRP \textless{} infinity when F\_WDM \textless{} 0.6 (frequentist 99.7\% CL) or F\_WDM \textless{} 0.35 (Bayesian 95\%). The pure-WDM bounds are m\_NRP \textgreater{} 8 keV (99.7\%) and 12.1 keV (Bayesian 95\%).

\textbf{Finding:} NuSTAR Galactic Center observations (Perez+2017, arXiv:1609.00667) give, by the abstract's own description, the strongest X-ray line limits on sterile-neutrino mixing over most of 10-50 keV in mass, improving earlier limits by up to about 10x. This is the published limit set that covers the DW-equivalent masses of 29-46 keV, which the server's tabulated limits do not cover.

\textbf{Finding:} I made a DW exclusion plot over m\_s = 1-100 keV and sin\textasciicircum{}2 2theta = 1e-14 to 1e-6, from an 80x80 scan (6400 points): 3092 points overclose and 968 are BBN-excluded. The plot shows the overclosure region, the Omega\_DM line, the NuSTAR/XMM limits, the Lyman-alpha thermal-WDM bound and the 3.5 keV point. It is Figure 10.

\textbf{Completed --- Build the X-ray decay and mixed cold-plus-warm side of the comparison. This means X-ray mixing-angle limits at the DW-equivalent masses, the predicted line energies and fluxes, and the published Lyman-alpha bounds on the warm fraction f versus mass..} Partly done: the X-ray limits above 26 keV and the DW-yield check (step 4) are still open. (1) The sterile-nu server does compute keV-scale DW relic densities (answers q-01). Its DW relic curve (Omega\_s h\textasciicircum{}2 = 0.12) gives sin\textasciicircum{}2 2theta = 2.545e-10, 3.946e-11, 2.022e-11, 1.084e-11 and 8.81e-12 at m\_s = 7.1, 20, 29, 41 and 46 keV, with radiative lifetimes of 1.6e27, 5.82e25, 1.77e25, 5.85e24 and 4.05e24 s. The tool flags 20, 29 and 41 keV as 'overcloses' even though Omega h\textasciicircum{}2 = 0.12, which looks like a rounding quirk. Results are in sterile\_nu\_grid\_DW\_e67bfc.csv in the run directory (finding f-001). (2) At those mixings the line energies are 3.55, 10, 14.5, 20.5 and 23 keV, and the whole-halo M31 fluxes are 1.75e-3, 1.71e-2, 3.88e-2, 8.31e-2 and 0.107 ph/s/cm\textasciicircum{}2. The tool's tabulated limits exclude pure DW at 7.1 keV (NuSTAR M31 Roach+2020 and XMM blank sky Dessert+2020) and at 20 keV (NuSTAR M31). It has no tabulated limit within 30\% in mass at 29, 41 or 46 keV, so those masses are uncovered rather than allowed (f-002). (3) The published limits covering 10-50 keV are from NuSTAR Galactic Center observations (Perez+2017, arXiv:1609.00667), which the abstract describes as the strongest over most of that range. Only the abstract was read, so numerical sin\textasciicircum{}2 2theta limits at 29-46 keV were not extracted (f-008, q-06). (4) For the mixed-model Lyman-alpha bound, Boyarsky+2009 (arXiv:0812.0010, cite-validated in references.bib) allow any m\_NRP \textgreater{} 5 keV when F\_WDM \textless{} 0.6 (99.7\% frequentist) or F\_WDM \textless{} 0.35 (Bayesian 95\%). Their pure-WDM bounds are m\_NRP \textgreater{} 8 keV (99.7\%) and 12.1 keV (95\%) (f-007). Pure DW is excluded by the tool's X-ray limits for m\_s up to about 20 keV. That rules out the conservative end of the Lyman-alpha DW-equivalent range, but exclusion at 29-46 keV still needs the NuSTAR numbers. Not done: the papers 1702.08430, 1705.01837, 1706.03118, 2007.13751, 2006.09553 and 1512.05369 were not read, so no DW-yield revision was checked. No figures were produced.

\textbf{Completed --- Use the sterile-nu server to compute non-resonant (Dodelson-Widrow) production fractions Omega\_s/Omega\_dm and related observables across keV masses and mixing angles. Then locate the DW relic-abundance curve (Omega\_s h\textasciicircum{}2 = 0.12) in the (m\_s, sin\textasciicircum{}2 2theta) plane..} Using the sterile-nu server (serve-sbbn), I ran a DW (non-resonant) grid of 15 masses from 1 to 50 keV times 13 mixings from 1e-13 to 1e-7 (CSV sterile\_nu\_grid\_DW\_8d35ed.csv), plus a follow-up grid at the five DW-equivalent masses (sterile\_nu\_grid\_DW\_10ad91.csv). On this server Omega\_s h\textasciicircum{}2 is exactly linear in sin\textasciicircum{}2 2theta and grows roughly as m\_s\textasciicircum{}1.8. Setting Omega\_s h\textasciicircum{}2 = 0.12 gives the DW relic curve: sin\textasciicircum{}2 2theta = 8.7e-9 (1 keV), 1.2e-9 (3 keV), 4.8e-10 (5 keV), 2.5e-10 (7.1 keV), 1.4e-10 (10 keV), 3.9e-11 (20 keV), 3.3e-11 (22 keV), 2.0e-11 (29 keV), 1.08e-11 (41 keV), 8.8e-12 (46 keV), 7.6e-12 (50 keV) (finding f-003). Fractions f = Omega\_s/Omega\_dm: 0.196 at 7.1 keV with sin\textasciicircum{}2 2theta = 5e-11. At sin\textasciicircum{}2 2theta = 1e-12: 0.025 (20 keV), 0.030 (22 keV), 0.049 (29 keV), 0.092 (41 keV), 0.114 (46 keV). These are ten times larger at 1e-11, where 46 and 50 keV overclose (finding f-004). Free-streaming length: 151 kpc at 7.1 keV, 107 kpc at 10 keV, 54 kpc at 20 keV. Delta N\_eff is at most about 3e-4, so BBN does not constrain DW keV dark matter: the BBN-allowed mixing is 1.5e-8 at about 50 keV, far above the relic curve. The server reports no \textless{}p/T\textgreater{} or momentum distribution and no per-point Lyman-alpha flag, so the \textless{}p/T\textgreater{} question cannot be answered with this tool. The scan scan\_DW\_d94a06.npz covers 1-100 keV and sin\textasciicircum{}2 2theta from 1e-14 to 1e-6 (6400 points; 3092 overclose, 968 BBN-excluded). The exclusion plot is Figure 10: /Users/nramachandra/hep-genesis-testing/projects/ramp-combo-lyman-wdm/runs/261019\_143114/exclusion\_scan\_DW\_d94a06.png (+ .pdf) (finding f-009). The X-ray-allowed f at each mass is (X-ray limit / relic-curve mixing) and needs the X-ray limits from d-03.

\textbf{Completed --- Verify the recent Lyman-alpha thermal-relic WDM limits and their systematic variations from the papers' text, and verify the thermal-relic to DW sterile-neutrino mass conversion..} Conversion checked against the paper text. Viel+2005 (astro-ph/0501562), Eq. 5, gives m\_s(DW) = 4.43 keV (m\_x/1 keV)\textasciicircum{}\{4/3\} (0.25*0.7\textasciicircum{}2/omega\_x)\textasciicircum{}\{1/3\}, with the omega\_x normalization at 0.1225 (tc-...-31; cited in references.bib, tc-...-41). The paper says the masses map one-to-one when the transfer functions match, which assumes a rescaled Fermi-Dirac spectrum. The survey's 4.46 keV prefactor does not appear in the paper. It equals 4.43 keV rescaled to omega\_dm = 0.12, since 4.43*(0.1225/0.12)\textasciicircum{}\{1/3\} = 4.461 keV. I worked out the DW-equivalent masses by hand, not with a tool, at omega\_dm = 0.1225 / 0.12: m\_x = 3.1 keV gives 20.0/20.2 keV; 3.3 gives 21.8/21.9; 3.5 gives 23.5/23.7; 4.1 gives 29.1/29.3; 4.5 gives 32.9/33.1; 5.3 gives 40.9/41.2; 5.7 gives 45.1/45.4 keV. The survey's 20-46 keV holds to within 0.6 keV (finding f-005).

Verified limits (finding f-006):

\begin{itemize}
  \item Viel+2013 (1306.2314, HIRES/MIKE, z\textgreater{}4): m\_x \textgreater{} 3.3 keV (2 sigma).
  \item Irsic+2017 (1702.01764): XQ-100 alone 1.4 keV, HIRES/MIKE 4.1 keV, combined 5.3 keV (2 sigma). Allowing temperature jumps of up to 5000 K lowers the combined limit to 3.5 keV.
  \item Palanque-Delabrouille+2019 (1911.09073, SDSS DR14 z\textless{}4.5 plus XQ-100): 5.3 keV (95\% CL).
  \item Villasenor+2022 (2209.14220, 1080 Cholla runs, Boera+2019 data): 3.1 keV (95\% CL), likelihood peak at 4.5 keV. The suppression saturates at k \textgreater{} 0.1 s/km because of peculiar-velocity smearing.
  \item Irsic+2023 (2309.04533, HIRES/UVES z=4.2-5.0, k\_max = 0.2 s/km): 5.7 keV (95\% CL), or 4.1 keV with k\_max \textless{} 0.1 s/km. Peculiar velocities and inhomogeneous reionization are the main small-scale systematics.
\end{itemize}
Verbatim quotes were cited only for astro-ph/0501562, 2309.04533 and 1702.01764, in /Users/nramachandra/hep-genesis-testing/projects/ramp-combo-lyman-wdm/runs/261019\_143114/references.bib. The other three limits come from arXiv abstracts (get\_arxiv\_metadata) and have not been cited yet. No figures were extracted.

Still open (question q-05): no post-2023 limit was found. The DESI search returned only BAO papers (dead end x-01). It is also unchecked whether 2209.14220, 1911.09073 or 1306.2314 quote their own DW or non-resonant sterile-neutrino masses or an updated conversion such as a Bozek/Abazajian-type relation.

\textbf{Critic, wave 1} --- continue: 7 findings stand up to the evidence, 2 questioned ('At the DW relic mixing, the line energies are E = 3.55, 10, 14.5, 20.5 and 23 keV for m\_s ': The line energies, lifetimes and M31 fluxes are visible in tc-36 to tc-40. The central claim is not: none of the quoted excerpts shows the tool's tabulated limits or the exclusion of m\_s = 7.1 keV (Roach+2020, Dessert+2020) and 20 keV (NuSTAR M31). All five excerpts are cut off before any limit comparison. The exclusion statement needs an excerpt that shows the limit values and the comparison verdict.; 'DW fractions f = Omega\_s h\textasciicircum{}2 / 0.12 from the server: at 7.1 keV and sin\textasciicircum{}2 2theta = 5e-11, ': The claim that Delta N\_eff is at most 3e-4 near the DW relic curve is contradicted by the board's own evidence. tc-261019\_1-w1-d-03-35 gives delta\_neff = 0.001574 at m\_s = 7.1 keV on the relic curve (sin\textasciicircum{}2 2theta = 2.545e-10, Omega\_s h\textasciicircum{}2 = 0.12). The 3e-4 value (tc-24) is for sin\textasciicircum{}2 2theta = 5e-11, which is 20\% of the relic abundance. BBN is still unconstraining, but the stated bound is about 5x too small. The fractions f = 0.196 (7.1 keV) and 0.030 (22 keV) do check out against tc-24 and tc-47.).

Issues: The cross-check against a published DW yield requested in c-01 was not done, and the server normalization looks off. The Abazajian (2005) fit, Omega\_s h\textasciicircum{}2 \textasciitilde{} 0.1 (sin\textasciicircum{}2 2theta/3e-9)\textasciicircum{}1.23 (m\_s/3 keV)\textasciicircum{}1.8, needs about 3e-9 at 3 keV for the full relic abundance. The server gives 1.2e-9, a factor of about 2.5 more efficient. The server is also exactly linear in sin\textasciicircum{}2 2theta, where the fit has power 1.23. The literature result is that DW at the 3.5 keV-line mixing (about 7e-11) yields only about 1\% of the dark matter. The server's f = 0.196 at 5e-11 is roughly 10-20x higher. The direction of the overall conclusion likely survives: a lower DW yield means larger relic mixing and stronger X-ray exclusion. The tabulated f values and relic-curve mixings should not be reported as literature-grade until this is reconciled. Note also that the arithmetic itself checks out: 1 keV gives 0.12/1.384e-5 x 1e-12 = 8.67e-9, 22 keV gives 3.32e-11, and the slope is 22\textasciicircum{}1.8 = 261.; The lifetime values in f-001 are not in its cited evidence (tc-26, tc-35); they appear only in tc-36 to tc-40. In addition, 3 of the 5 points at Omega\_s h\textasciicircum{}2 = 0.12 are flagged 'overcloses', which shows the overclosure threshold is applied with a strict or rounded inequality. Points sitting exactly on the relic curve should not be labelled overclosing in any table or figure.; The verdict on pure DW at 29-46 keV (the Lyman-alpha-allowed masses under the stronger limits) still rests on NuSTAR numbers that have not been extracted (q-06). Without them, the headline 'pure DW excluded' is shown only for the conservative 20 keV end. tc-51 (a two-sterile-neutrino oscillation paper) is irrelevant evidence for f-008..

\textbf{Next steps after wave 1} --- All earlier directions are closed and none are carried over, so this final wave takes the critic's three suggested follow-ups, trimmed to fit 14 tool calls (about 5, 5 and 4). The highest priority is closing the comparison itself: numerical NuSTAR limits at 29-46 keV and the X-ray-allowed fraction against the Boyarsky mixed bound (q-04, q-06), plus fixing contested f-002. Equal priority goes to the robustness cross-check of the server's DW yield normalization, which could shift f by an order of magnitude, and to correcting the contested Delta N\_eff claim in f-004. Lower priority is the mapping-robustness and post-2023 literature check (q-02, q-05), which tests whether the 20-46 keV range is stable. The task does not request facility runs, so no HPC dispatch is needed. Figure 10's limit-set provenance should be documented in the final write-up rather than through a new direction.

New lines of investigation: Extract numerical NuSTAR X-ray upper limits on sin\textasciicircum{}2 2theta at m\_s = 20, 29, 41 and 46 keV. Combine them with the DW relic curve to get the X-ray-allowed fraction f\_X(m\_s), compare f\_X with the Boyarsky+2009 mixed cold-plus-warm bound, and give a final verdict on pure non-resonant production over 20-46 keV.; Cross-check the sterile-nu server's DW yield normalization and its sin\textasciicircum{}2 2theta scaling against published production calculations. Derive a correction factor and propagate it to the relic curve (f-003), the fractions f (f-004) and the Delta N\_eff statement.; Test how robust the thermal-relic to NRP mass mapping is by tabulating paper-quoted NRP-equivalent masses next to the Viel+2005-derived values. Run one targeted search for post-2023 Lyman-alpha WDM limits from DESI P1D or emulator re-analyses..

\textbf{Wave 2} --- now exploring, in parallel: Extract numerical NuSTAR X-ray upper limits on sin\textasciicircum{}2 2theta at m\_s = 20, 29, 41 and 46 keV. Combine them with the DW relic curve to get the X-ray-allowed fraction f\_X(m\_s), compare f\_X with the Boyarsky+2009 mixed cold-plus-warm bound, and give a final verdict on pure non-resonant production over 20-46 keV.; Cross-check the sterile-nu server's DW yield normalization and its sin\textasciicircum{}2 2theta scaling against published production calculations. Derive a correction factor and propagate it to the relic curve (f-003), the fractions f (f-004) and the Delta N\_eff statement.; Test how robust the thermal-relic to NRP mass mapping is by tabulating paper-quoted NRP-equivalent masses next to the Viel+2005-derived values. Run one targeted search for post-2023 Lyman-alpha WDM limits from DESI P1D or emulator re-analyses..

\textbf{Finding:} Comparing mappings from thermal-relic mass to NRP/DW mass (my own arithmetic, no tool used): Viel+2005 Eq. 5 (4.43 keV m\_x\textasciicircum{}\{4/3\}, omega=0.1225) gives 20.0, 21.8, 23.5, 29.1, 40.9 and 45.1 keV for m\_x = 3.1, 3.3, 3.5, 4.1, 5.3 and 5.7 keV. A Bozek+2016-type relation, 3.9 keV m\_x\textasciicircum{}\{1.294\}, gives 16.9, 18.3, 19.7, 24.2, 33.8 and 37.1 keV. That is lower by 3.1, 3.5, 3.8, 4.9, 7.1 and 8.0 keV, so the DW-equivalent range moves from 20-45 keV to 17-37 keV. The 3.9 keV and 1.294 coefficients come from the direction card. I could not confirm them from any paper text: the abstract search found nothing, and the 2007.13751 abstract quotes no relation.

\textbf{Finding:} Of the abstracts checked, only Boyarsky+2009 (0812.0010) quotes non-resonant (NRP) sterile-neutrino masses directly: m\_NRP \textgreater{} 8 keV (99.7\%, frequentist) and \textgreater{} 12.1 keV (Bayesian 95\%). Inverting Viel+2005 by hand, these correspond to thermal-relic masses of about 1.56 and 2.12 keV, which are older and weaker than the 3.1-5.7 keV limits. The abstracts of Irsic+2017 (1702.01764) and Palanque-Delabrouille+2019 (1911.09073) give only thermal-relic masses (5.3 keV for both; 3.5 keV for Irsic+2017 with temperature jumps), so the DW-equivalent values for them come entirely from a mapping relation.

\textbf{Finding:} Perez+2017 (arXiv:1609.00667, NuSTAR Galactic Center 0-bounce photons) states in the text that, within the nuMSM, its limit gives an upper bound on the sterile-neutrino mass of m\_chi below about 16 keV. The nuMSM allowed band in its Fig. 8 lies between the non-resonant (DW) production line above and the maximum-lepton-asymmetry resonant line below. So the NuSTAR GC limit lies below the whole band, including the DW line, at masses above about 16 keV and up to the edge of its 10-50 keV coverage. Above 50 keV, INTEGRAL limits rule out mixing-produced sterile neutrinos as all of the dark matter. Together these imply pure DW is X-ray excluded at m\_s = 20, 29, 41 and 46 keV. This is an inference from the paper's text, not from numerical sin\textasciicircum{}2 2theta values. The paper notes the limit weakens near 13 and 40 keV ('bumps'), which is a caveat for 41 keV.

\textbf{Finding:} Perez+2017 Eq. 5 gives the GC decay-line flux as F $\approx$ 2.6e-6 cm\textasciicircum{}-2 s\textasciicircum{}-1 (m/20 keV)\textasciicircum{}4 (sin\textasciicircum{}2 2theta/1e-14)(dOmega/4 deg\textasciicircum{}2)(J/40 GeV cm\textasciicircum{}-3 kpc sr\textasciicircum{}-1), with dOmega\_tot $\approx$ 3.8 deg\textasciicircum{}2 and J\_tot = 29 GeV cm\textasciicircum{}-3 kpc sr\textasciicircum{}-1 for the default coreNFW profile. Applied to the DW relic mixing at 20 keV (3.9e-11), this predicts about 7e-3 ph/cm\textasciicircum{}2/s in the NuSTAR 0-bounce aperture (my hand calculation). Numerical flux and sin\textasciicircum{}2 2theta limits are shown only in Figs. 6-8, which are not machine-readable here.

\textbf{Finding:} The sterile-nu tool's X-ray verdicts at the DW relic mixing give only a list of excluding limits, never the limit value itself. At m\_s = 20 keV, sin\textasciicircum{}2 2theta = 3.946e-11, the M31 flux is 0.0171 ph/s/cm\textasciicircum{}2 and the result is 'EXCLUDED by: NuSTAR M31 (Roach+ 2020)'. At m\_s = 29 keV, sin\textasciicircum{}2 2theta = 2.022e-11 (E = 14.5 keV, tau = 1.77e25 s, M31 flux 0.0388, MW-center whole-halo flux 264 ph/s/cm\textasciicircum{}2), the result is 'Not excluded by tabulated X-ray limits (within 30\% in mass)' with excluded\_by = []. So f-002 cannot be backed by a numerical limit from this tool, and its 29 keV 'not excluded' is a coverage gap, not an allowed region.

\textbf{Completed --- Test how robust the thermal-relic to NRP mass mapping is by tabulating paper-quoted NRP-equivalent masses next to the Viel+2005-derived values. Run one targeted search for post-2023 Lyman-alpha WDM limits from DESI P1D or emulator re-analyses..} I built a side-by-side table of NRP/DW-equivalent masses under two relations (finding f-010). I did the arithmetic by hand; no tool computed it. The relations are Viel+2005 Eq. 5 (4.43 keV m\_x\textasciicircum{}\{4/3\}, omega = 0.1225; verified in d-02) and a Bozek+2016-type relation (3.9 keV m\_x\textasciicircum{}\{1.294\}). I took the Bozek coefficients from the direction card and could not confirm them in any paper text: the abstract search found nothing and the 2007.13751 abstract quotes no relation. Treat that column as unverified.

\begin{longtable}{@{}p{0.320\textwidth}p{0.320\textwidth}p{0.320\textwidth}@{}}
\toprule
\textbf{m\_x (keV)} & \textbf{Viel+2005 (keV)} & \textbf{Bozek-type (keV)} \\ \midrule \endhead
3.1 & 20.0 & 16.9 \\
3.3 & 21.8 & 18.3 \\
3.5 & 23.5 & 19.7 \\
4.1 & 29.1 & 24.2 \\
5.3 & 40.9 & 33.8 \\
5.7 & 45.1 & 37.1 \\
\bottomrule\end{longtable}
Under the alternative relation the DW-equivalent range moves from 20-45 keV to 17-37 keV, lower by 3-8 keV.

Of the abstracts checked, only Boyarsky+2009 (0812.0010) quotes NRP masses itself: m\_NRP \textgreater{} 8 keV (99.7\%) and \textgreater{} 12.1 keV (Bayesian 95\%). Inverting Viel+2005 by hand gives m\_x of about 1.56 and 2.12 keV, weaker than modern limits. The Irsic+2017 (1702.01764) and Palanque-Delabrouille+2019 (1911.09073) abstracts quote only thermal-relic masses (5.3 keV; 3.5 keV for Irsic+2017 with temperature jumps), so their DW masses depend entirely on a mapping relation (finding f-011).

The post-2023 search found nothing (dead end x-02). Neither query ('one-dimensional flux power spectrum warm dark matter 2024' or 'Lyman-alpha P1D emulator thermal relic warm dark matter mass limit DESI eBOSS') returned a DESI DR1, eBOSS or emulator WDM limit, only DESI BAO papers (2404.03001, 2503.14739).

Still open: q-05 should be retried with web\_search or arxiv\_agent using author names; the Bozek+2016 relation needs checking against its source; and I did not check whether 2209.14220 and 1306.2314 quote their own NRP masses. No files or figures were produced.

\textbf{Blocked --- Extract numerical NuSTAR X-ray upper limits on sin\textasciicircum{}2 2theta at m\_s = 20, 29, 41 and 46 keV. Combine them with the DW relic curve to get the X-ray-allowed fraction f\_X(m\_s), compare f\_X with the Boyarsky+2009 mixed cold-plus-warm bound, and give a final verdict on pure non-resonant production over 20-46 keV..} (The limit curves exist only in figures of the paper's PDF that the figure tool could not extract, and the X-ray tool reports a verdict without a limit value. Numerical f\_X therefore needs the PDF figures or the Roach+2020/2023 data, and the tool budget ran out before I could get them.) I could not get numerical NuSTAR sin\textasciicircum{}2 2theta limits at 20, 29, 41 or 46 keV, so there is no f\_X(m\_s) table and no f\_X comparison with the Boyarsky+2009 bound (F\_WDM \textless{} 0.35-0.6). Pure DW comes out excluded across 20-46 keV, but only by inference from the Perez+2017 text, not from numbers.

What was established: (1) Perez+2017 (arXiv:1609.00667; full text in run\_dir/1609.00667.txt, tc-...-69, -76, -78). The paper states 'Assuming the $\nu$MSM framework, our limit translates into an upper limit on the sterile neutrino mass of m$\chi$ ? 16 keV'. In its Fig. 8 the upper edge of the nuMSM band is the non-resonant (DW) production line. So the NuSTAR GC limit lies below the DW relic mixing for m\_s above about 16 keV, across its 10-50 keV coverage. INTEGRAL covers m\_s above 50 keV and is said to rule out mixing-produced sterile neutrinos as all of the dark matter. My inferred verdict is that pure DW is X-ray excluded at 20, 29, 41 and 46 keV. Caveat: the paper notes a weakening ('bump') of its limit near 40 keV, so 41 keV is the least secure point (f-012, medium confidence). (2) Perez+2017 Eq. 5 gives the flux normalization F $\approx$ 2.6e-6 cm\textasciicircum{}-2 s\textasciicircum{}-1 (m/20 keV)\textasciicircum{}4 (sin\textasciicircum{}2 2theta/1e-14) for the default coreNFW profile (J\_tot = 29 GeV cm\textasciicircum{}-3 kpc sr\textasciicircum{}-1, dOmega $\approx$ 3.8 deg\textasciicircum{}2). From this, the DW relic point at 20 keV predicts about 7e-3 ph/cm\textasciicircum{}2/s; this is my hand calculation (f-013). (3) The sterile-nu tool returns no limit values. At 20 keV and sin\textasciicircum{}2 2theta = 3.946e-11 it reports EXCLUDED by NuSTAR M31 (Roach+2020), with an M31 flux of 0.0171 ph/s/cm\textasciicircum{}2. At 29 keV and 2.022e-11 it reports excluded\_by = [] ('not excluded within 30\% in mass'), with E = 14.5 keV and tau = 1.77e25 s. That 29 keV result is a coverage gap, not an allowed region. f-002 therefore cannot be backed by a quoted tool limit value; it rests on the verdict strings only (f-014).

What failed: get\_paper\_figures could not render Figs. 6-8 of 1609.00667, since the HTML has only Figs. 3 and 4 (dead end x-03). Roach+2020/2023 were not checked, and Perez+2017 was not added to references.bib because the budget ran out.

Still open: q-06 and q-07 (numerical limits at 29-46 keV and the f\_X table), and q-04 (comparing f\_X with Boyarsky+2009).

\textbf{Finding:} Abazajian 2005 (astro-ph/0511630, Eq. 7, full numerical DW calculation with the QCD transition, nu\_e mixing): m\_s = 3.40 keV (sin\textasciicircum{}2 2theta/1e-8)\textasciicircum{}(-0.615) (Omega\_DM/0.26)\textasciicircum{}0.5 x \{0.527 erfc[-1.15 (T\_QCD/170 MeV)\textasciicircum{}2.15]\}. The braced factor is 1 at T\_QCD = 170 MeV, and the fit is valid for 135 \textless{} T\_QCD \textless{} 300 MeV. Inverted, this gives Omega\_s = 0.26 (m\_s/3.4 keV)\textasciicircum{}2 (sin\textasciicircum{}2 2theta/1e-8)\textasciicircum{}1.23. Compared with the server's linear scaling in mixing and roughly m\textasciicircum{}1.8 in mass, the literature yield goes as mixing\textasciicircum{}1.23 and m\_s\textasciicircum{}2. The peak production temperature is T\_peak \textasciitilde{} 130 MeV (m\_s/3 keV)\textasciicircum{}(1/3).

\textbf{Finding:} Literature DW relic mixing from the Abazajian 2005 Eq. 7 (my hand calculation, Omega\_DM = 0.26, T\_QCD = 170 MeV): sin\textasciicircum{}2 2theta = 1.23e-8 (3 keV), 3.02e-9 (7.1 keV), 5.61e-10 (20 keV), 3.06e-10 (29 keV), 1.75e-10 (41 keV), 1.45e-10 (46 keV). The server's relic curve (f-003) lies 10.2x, 12.1x, 14.4x, 15.3x, 16.1x and 16.4x lower in mixing at these masses. So the server overproduces DW sterile neutrinos at fixed mixing by more than the critic's 2.5x estimate. The f-001 and f-003 relic mixings should be multiplied by about 10-16 to match this fit. Using Omega\_DM = 0.264 (h = 0.674) raises the literature mixings by only about 1\%.

\textbf{Finding:} Revised DW fractions under the Abazajian 2005 yield, f = (sin\textasciicircum{}2 2theta / sin\textasciicircum{}2 2theta\_relic)\textasciicircum{}1.23 normalized to Omega\_DM = 0.26 (my hand calculation): at 7.1 keV, f = 0.0064 for sin\textasciicircum{}2 2theta = 5e-11 and 0.0098 for 7e-11. This matches the literature statement of about 1\% of the dark matter at the 3.5 keV-line mixing, whereas the server gives 0.196 and 0.27 (about 30x higher). At sin\textasciicircum{}2 2theta = 1e-12: f = 4.2e-4 (20 keV), 8.8e-4 (29 keV), 1.75e-3 (41 keV), 2.2e-3 (46 keV), against the server's 0.025, 0.049, 0.092 and 0.114 (50-60x too high). At 1e-11, multiply by 10\textasciicircum{}1.23 = 17: f = 7e-3 (20 keV) to 0.037 (46 keV), so nothing overcloses. Overall, f-004's fractions are overstated by 1.5-2 orders of magnitude.

\textbf{Finding:} Corrected Delta N\_eff statement: on the relic curve the sterile number density is fixed by Omega\_s h\textasciicircum{}2 / m\_s, whatever yield normalization is used. My hand estimate, assuming one Weyl species and n\_s/n\_nu = Omega\_s h\textasciicircum{}2 x 94 eV / m\_s, gives Delta N\_eff $\approx$ 11.3 eV / m\_s while relativistic: 1.6e-3 at 7.1 keV, 5.6e-4 at 20 keV, 2.5e-4 at 46 keV. This matches the critic's 1.6e-3 on-curve value at 7.1 keV, so f-004's 'at most 3e-4' holds only above about 40 keV. All values are 2-3 orders of magnitude below BBN and Planck sensitivity (about 0.2-0.3), so BBN does not constrain keV DW dark matter.

\textbf{Completed --- Cross-check the sterile-nu server's DW yield normalization and its sin\textasciicircum{}2 2theta scaling against published production calculations. Derive a correction factor and propagate it to the relic curve (f-003), the fractions f (f-004) and the Delta N\_eff statement..} Two success criteria are met (relic-mixing ratios at six masses, revised fractions f with the normalization stated) and the corrected Delta N\_eff statement is recorded as an analytic estimate. The cite() step was not done. All numbers below except the fit itself are my hand arithmetic from that fit, not tool output.

(1) Literature yield (f-015). Abazajian 2005 (astro-ph/0511630, Eq. 7), read verbatim, gives m\_s = 3.40 keV (sin\textasciicircum{}2 2theta/1e-8)\textasciicircum{}-0.615 (Omega\_DM/0.26)\textasciicircum{}0.5 x \{0.527 erfc[-1.15 (T\_QCD/170 MeV)\textasciicircum{}2.15]\}. The braced factor equals 1 at T\_QCD = 170 MeV, and the fit is valid for 135 \textless{} T\_QCD \textless{} 300 MeV. This means Omega\_s scales as (sin\textasciicircum{}2 2theta)\textasciicircum{}1.23 m\_s\textasciicircum{}2, whereas the server is linear in mixing and goes roughly as m\_s\textasciicircum{}1.8. Paper text: astro-ph\_0511630.txt in the run directory.

(2) Relic-curve correction (f-016). The literature relic mixings are 1.23e-8 (3 keV), 3.02e-9 (7.1 keV), 5.6e-10 (20 keV), 3.06e-10 (29 keV), 1.75e-10 (41 keV) and 1.45e-10 (46 keV). These are 10.2x, 12.1x, 14.4x, 15.3x, 16.1x and 16.4x the server's f-003 values, so the critic's 2.5x estimate understates the gap. The f-001 and f-003 mixings should be scaled up by about 10-16.

(3) Revised fractions (f-017). At 7.1 keV, f = 0.0064 for sin\textasciicircum{}2 2theta = 5e-11 and 0.0098 for 7e-11, consistent with the literature figure of about 1\%. The server gives 0.196 at 5e-11. At sin\textasciicircum{}2 2theta = 1e-12, f = 4.2e-4 (20 keV), 8.8e-4 (29 keV), 1.75e-3 (41 keV) and 2.2e-3 (46 keV), 50-60x below f-004. At 1e-11 nothing overcloses.

(4) Delta N\_eff (f-018). On the relic curve Delta N\_eff $\approx$ 11.3 eV / m\_s: 1.6e-3 at 7.1 keV, 5.6e-4 at 20 keV and 2.5e-4 at 46 keV. This confirms the critic's 1.6e-3, and f-004's 'at most 3e-4' holds only above about 40 keV. BBN and the CMB do not constrain these values. This is an analytic estimate that does not depend on the yield normalization.

Caveats: Asaka-Laine-Shaposhnikov (hep-ph/0612182, abstract only) find results that 'differ moderately' and allow DW as all of the dark matter only below 3.5 keV (6 keV at most). The updated yields in 1705.01837 and Venumadhav+2016 were not read.

Still open (q-08): the cite() call with the verbatim Eq. 7 quote into references.bib, and checking the 2017 and 2016 yields.

\textbf{Critic, wave 2} --- continue: 3 findings stand up to the evidence, 6 questioned ('Comparing mappings from thermal-relic mass to NRP/DW mass (my own arithmetic, no tool used': The arithmetic is correct: 3.9*3.1\textasciicircum{}1.294 = 16.9 and 3.9*5.7\textasciicircum{}1.294 = 37.1. However, none of the evidence excerpts contains the 3.9 keV / 1.294 relation. tc-82 is an unrelated search hit about UHE neutrino damping and two-sterile oscillations. tc-84 is an abstract that, as the finding itself admits, quotes no relation. The claimed shift of the DW range to 17-37 keV therefore rests on unverified coefficients. This relation is commonly attributed to Abazajian 2006 (astro-ph/0512631), not Bozek+2016, and should be checked against that source.; 'Perez+2017 (arXiv:1609.00667, NuSTAR Galactic Center 0-bounce photons) states in the text ': Neither quoted excerpt contains the claimed statements. tc-76 is the title and author block. tc-78 is the NFW-profile and local-density passage. The excerpts include neither 'm\_chi \textless{}\textasciitilde{} 16 keV', nor the Fig. 8 band-edge description, nor the 13/40 keV 'bump' remark, nor the INTEGRAL statement. The title supports the general thrust, but the specific claims are unevidenced. The exclusion at 29-46 keV is an inference from figures that were never read.; 'Perez+2017 Eq. 5 gives the GC decay-line flux as F $\approx$ 2.6e-6 cm\textasciicircum{}-2 s\textasciicircum{}-1 (m/20 keV)\textasciicircum{}4 (sin\textasciicircum{}2': Eq. 5, dOmega = 3.8 deg\textasciicircum{}2 and J\_tot = 29 do not appear in the quoted excerpt, which is the NFW-profile passage. The hand arithmetic is correct (2.6e-6 \textit{ 3900 } 0.95 \textit{ 0.725 = 7.0e-3). However, it uses the server relic mixing of 3.9e-11, which f-016 of this same wave says is about 14x too low. With the literature mixing of 5.6e-10, the predicted flux is about 0.1 ph/cm\textasciicircum{}2/s, so the finding is internally inconsistent with the wave's own correction.; 'Abazajian 2005 (astro-ph/0511630, Eq. 7, full numerical DW calculation with the QCD transi': The only content excerpt (tc-81) shows Eq. 4, the lepton-asymmetry potential, not Eq. 7, and tc-70 is a download path. The quoted Eq. 7 form matches the commonly cited Abazajian fit, and the inversion is algebraically correct, but the verbatim fit is absent from the evidence and the cite() step was skipped. The interpretation that the yield goes as mixing\textasciicircum{}1.23 is also physically misleading. At fixed mass, DW production is linear in sin\textasciicircum{}2 2theta for small mixing. The -0.615 exponent of the relic-curve fit encodes the mass dependence (g}(T\_peak)) along the Omega = const curve, not a nonlinear dependence on mixing.).

Issues: Recompute the fractions with f linear in mixing at fixed mass, using the literature relic mixing. The 1.23 exponent biases f low by a factor of 2.6-4.; There is still no numerical X-ray limit at 29-46 keV. The corrected literature DW relic mixing is 10-16x larger, which makes X-ray exclusion of pure DW much easier: the predicted 20 keV GC flux rises from about 7e-3 to about 0.1 ph/cm\textasciicircum{}2/s. The verdict should be restated against the corrected relic curve.; The tool labels its limit 'NuSTAR M31 (Roach+ 2020)'. Roach+2020 (1908.09037) is a Galactic-bulge analysis, while the NuSTAR M31 limit is Ng+2019 (1901.01262), so the provenance needs checking. The 29 keV excerpt (tc-72) is also truncated before the 'not excluded' verdict string..

Stopping after wave 2: budget reached (2/2 waves used). Writing the report from the board.

\textit{Research report:}\par
\subsubsection*{Lyman-alpha forest limits on warm dark matter and what they imply for non-resonant (Dodelson-Widrow) sterile neutrinos}
\subsubsection*{Abstract}
This report collects recent Lyman-alpha forest lower limits on the thermal-relic warm dark matter (WDM) mass and checks them against the papers. It converts those limits to non-resonantly produced (Dodelson-Widrow, DW) sterile-neutrino masses and compares the result with DW production yields and X-ray decay limits.

\textbf{Lyman-alpha limits.} The verified 95\% CL (2 sigma) limits run from m\_x \textgreater{} 3.1 keV (Villasenor+2022) to m\_x \textgreater{} 5.7 keV (Irsic+2023). Choices of IGM thermal history and k\_max weaken them to 3.5--4.1 keV.

\textbf{Mass mapping.} The Viel+2005 relation was checked against its text and maps these limits to m\_s(DW) $\approx$ 20--45 keV. An alternative Bozek-type relation, whose source was not confirmed, gives 17--37 keV.

\textbf{Production yields.} A sterile-neutrino production server was available and was used to compute DW relic densities and fractions. Its yield normalization disagrees with the Abazajian 2005 fit by a factor of 10--16 in relic mixing angle. The server-based fractions are therefore overstated by 1.5--2 orders of magnitude.

\textbf{X-ray side.} The server's tabulated X-ray limits exclude pure DW dark matter at 7.1 and 20 keV and have no coverage at 29--46 keV. The text of the NuSTAR Galactic Center analysis (Perez+2017) implies exclusion above about 16 keV. This conclusion is reached by inference from the paper's wording, without numerical mixing-angle limits.

\textbf{Overall.} Pure DW sterile neutrinos as all of the dark matter appear excluded: Lyman-alpha forces m\_s ? 17--20 keV, and X-ray data exclude DW above about 16 keV. The quantitative X-ray-allowed fraction could not be tabulated.

\subsubsection*{Background}
\paragraph{Thermal-relic limits from the Lyman-alpha flux power spectrum}
Flux-power-spectrum analyses constrain the free-streaming of WDM, usually quoted as a lower limit on the mass m\_x of a thermal relic. The papers covered here are:

\begin{itemize}
  \item \textbf{Viel+2013} (arXiv:1306.2314): m\_x \textgreater{} 3.3 keV (2 sigma) from HIRES/MIKE spectra at z \textgreater{} 4.
  \item \textbf{Irsic+2017} (arXiv:1702.01764): combines XQ-100 with HIRES/MIKE.
  \item \textbf{Palanque-Delabrouille+2019} (arXiv:1911.09073): uses SDSS DR14 BOSS/eBOSS together with XQ-100.
  \item \textbf{Villasenor+2022} (arXiv:2209.14220): an independent suite of 1080 Cholla simulations applied to the Boera+2019 data.
  \item \textbf{Irsic+2023} (arXiv:2309.04533): pushes to smaller scales with HIRES/UVES data at z = 4.2--5.0.
\end{itemize}
\paragraph{From thermal-relic mass to sterile-neutrino mass}
For non-resonantly produced sterile neutrinos, the standard mapping from m\_x to m\_s comes from Viel+2005 (arXiv:astro-ph/0501562). It assumes the DW momentum spectrum is a rescaled Fermi-Dirac distribution.

Two papers provide alternatives to that mapping:

\begin{itemize}
  \item \textbf{Boyarsky+2009} (arXiv:0812.0010) give Lyman-alpha bounds for mixed cold-plus-warm models, expressed directly in non-resonant-production (NRP) masses.
  \item \textbf{Rogers \& Peiris} (arXiv:2007.13751) give an emulator framework that applies Lyman-alpha bounds to arbitrary transfer functions without going through the thermal-relic mapping.
\end{itemize}
\paragraph{Production and X-ray constraints}
On the production side, the following papers were consulted:

\begin{itemize}
  \item \textbf{Abazajian 2005} (arXiv:astro-ph/0511630) provides a fitting formula for the DW yield that includes the QCD transition.
  \item \textbf{Asaka, Laine \& Shaposhnikov} (arXiv:hep-ph/0612182) find moderately different yields.
  \item \textbf{arXiv:1512.05369} states that DW production cannot meet the relic abundance and the structure-formation bounds at the same time.
  \item The reviews arXiv:1702.08430 and arXiv:1705.01837, and arXiv:2006.09553 on thermalization, are relevant background but were not read.
\end{itemize}
On the X-ray side, NuSTAR Galactic Center observations (Perez+2017, arXiv:1609.00667) give, by the abstract's own description, the strongest line limits over most of the 10--50 keV mass range. They improve earlier limits by up to about 10x at some masses.

\paragraph{A note on the survey brief}
The brief stated that no sterile-neutrino production tools were available. In fact a sterile-neutrino server was reachable and was used. It provides DW relic density, Delta N\_eff, BBN, X-ray line and free-streaming outputs.

\subsubsection*{Results}
\paragraph{Verified Lyman-alpha thermal-relic limits}
The limits below were checked against paper text where full text was read (Irsic+2017, Irsic+2023) and against arXiv abstracts otherwise (Viel+2013, Palanque-Delabrouille+2019, Villasenor+2022).

\begin{longtable}{@{}p{0.240\textwidth}p{0.240\textwidth}p{0.240\textwidth}p{0.240\textwidth}@{}}
\toprule
\textbf{Analysis} & \textbf{Data} & \textbf{m\_x limit} & \textbf{Systematic variation} \\ \midrule \endhead
Viel+2013 (1306.2314) & HIRES/MIKE, z \textgreater{} 4 & \textgreater{} 3.3 keV (2 sigma) & 1, 2, 2.5 keV disfavoured at about 9, 4, 3 sigma \\
Irsic+2017 (1702.01764) & XQ-100 alone / HIRES/MIKE alone / combined & 1.4 / 4.1 / 5.3 keV (2 sigma) & Combined limit falls to 3.5 keV if temperature jumps of up to 5000 K are allowed \\
Palanque-Delabrouille+2019 (1911.09073) & SDSS DR14 z \textless{} 4.5 + XQ-100 & \textgreater{} 5.3 keV (95\% CL) & --- \\
Villasenor+2022 (2209.14220) & 1080 Cholla runs, Boera+2019 data, z = 4.0--5.2 & \textgreater{} 3.1 keV (95\% CL) & Likelihood peaks at 4.5 keV; flux-power suppression saturates at k ? 0.1 s/km because of peculiar-velocity smearing \\
Irsic+2023 (2309.04533) & HIRES/UVES, z = 4.2--5.0, k\_max = 0.2 s/km & \textgreater{} 5.7 keV (95\% CL) & 4.1 keV for k\_max \textless{} 0.1 s/km with earlier thermal priors; peculiar velocities and inhomogeneous reionization are the main small-scale systematics \\
\bottomrule\end{longtable}
The spread among the recent limits (3.1 to 5.7 keV) is comparable to the shifts from thermal-history and k\_max choices (5.3 to 3.5 keV, and 5.7 to 4.1 keV). The modelling of small scales is therefore the dominant uncertainty.

Three searches for post-2023 limits found nothing: one query targeting DESI, and two targeting the one-dimensional flux power spectrum (P1D), eBOSS and emulator re-analyses. All returned only DESI BAO papers or unrelated results.

\paragraph{Mapping thermal-relic limits to DW sterile-neutrino masses}
\textbf{The Viel+2005 relation.} Eq. 5 of Viel+2005, read in the paper text, is

m\_s(DW) = 4.43 keV (m\_x/1 keV)\textasciicircum{}\{4/3\} (0.25$\cdot$0.7$^2$/omega\_x)\textasciicircum{}\{1/3\},

normalized at omega\_x = 0.1225. The paper describes the correspondence as one-to-one when the transfer functions match.

The 4.46 keV prefactor quoted in the survey brief does not appear in the paper. It equals 4.43 keV rescaled to omega\_dm = 0.12, since 4.43 $\times$ (0.1225/0.12)\textasciicircum{}\{1/3\} = 4.461 keV. The brief's range of 20--46 keV holds to within 0.6 keV.

\textbf{Comparison with a Bozek-type relation.} The table also gives a Bozek+2016-type relation, 3.9 keV (m\_x/keV)\textasciicircum{}\{1.294\}. Its coefficients came from the planning notes and could not be confirmed in any paper text: a targeted abstract search found nothing, and the Rogers \& Peiris abstract quotes no such relation. All conversions below are hand arithmetic, not tool output. The critic marked the two-relation comparison as contested.

\begin{longtable}{@{}p{0.240\textwidth}p{0.240\textwidth}p{0.240\textwidth}p{0.240\textwidth}@{}}
\toprule
\textbf{m\_x (keV)} & \textbf{Viel+2005, omega = 0.1225 (keV)} & \textbf{Viel+2005, omega = 0.12 (keV)} & \textbf{Bozek-type, unverified (keV)} \\ \midrule \endhead
3.1 & 20.0 & 20.2 & 16.9 \\
3.3 & 21.8 & 21.9 & 18.3 \\
3.5 & 23.5 & 23.7 & 19.7 \\
4.1 & 29.1 & 29.3 & 24.2 \\
4.5 (Villasenor peak) & 32.9 & 33.1 & --- \\
5.3 & 40.9 & 41.2 & 33.8 \\
5.7 & 45.1 & 45.4 & 37.1 \\
\bottomrule\end{longtable}
The alternative relation lowers the DW-equivalent masses by 3.1--8.0 keV, moving the range from 20--45 keV to 17--37 keV.

\textbf{Papers that quote NRP masses directly.} Of the abstracts checked, only Boyarsky+2009 does so: m\_NRP \textgreater{} 8 keV (99.7\%, frequentist) and \textgreater{} 12.1 keV (Bayesian 95\%). Inverting Viel+2005 by hand, these correspond to m\_x $\approx$ 1.56 and 2.12 keV, which is weaker than the modern limits.

The Irsic+2017 and Palanque-Delabrouille+2019 abstracts quote only thermal-relic masses. Their DW-equivalent values therefore rest entirely on a mapping relation. It was not checked whether Villasenor+2022 or Viel+2013 quote their own NRP masses.

\paragraph{DW production on the sterile-neutrino server}
\textbf{Scans run.} Several separate runs were made on the server:

\begin{itemize}
  \item A 15-mass $\times$ 13-mixing DW grid spanning 1--50 keV and sin$^2$2$\theta$ = 1e-13 to 1e-7 (file \texttt{sterile\_\allowbreak{}nu\_\allowbreak{}grid\_\allowbreak{}DW\_\allowbreak{}8d35ed.csv}).
  \item A follow-up grid at 22, 29, 41, 46 and 50 keV with sin$^2$2$\theta$ = 1e-12 and 1e-11 (\texttt{sterile\_\allowbreak{}nu\_\allowbreak{}grid\_\allowbreak{}DW\_\allowbreak{}10ad91.csv}).
  \item A five-point run on the derived relic curve at 7.1, 20, 29, 41 and 46 keV (\texttt{sterile\_\allowbreak{}nu\_\allowbreak{}grid\_\allowbreak{}DW\_\allowbreak{}e67bfc.csv}).
  \item An 80 $\times$ 80 scan, described with Figure 10 below.
\end{itemize}
\textbf{Relic curve.} On this server, Omega\_s h$^2$ is exactly linear in sin$^2$2$\theta$ and grows roughly as m\_s\textasciicircum{}1.8. Setting Omega\_s h$^2$ = 0.12 gives the server's DW relic curve:

\texttt{\textbar{}\allowbreak{} m\_\allowbreak{}s (keV)\allowbreak{} \textbar{}\allowbreak{} 1 \textbar{}\allowbreak{} 2 \textbar{}\allowbreak{} 3 \textbar{}\allowbreak{} 5 \textbar{}\allowbreak{} 7.1 \textbar{}\allowbreak{} 10 \textbar{}\allowbreak{} 15 \textbar{}\allowbreak{} 20 \textbar{}\allowbreak{} 22 \textbar{}\allowbreak{} 29 \textbar{}\allowbreak{} 41 \textbar{}\allowbreak{} 46 \textbar{}\allowbreak{} 50 \textbar{}\allowbreak{}}\\
\texttt{\textbar{}\allowbreak{}---\textbar{}\allowbreak{}---\textbar{}\allowbreak{}---\textbar{}\allowbreak{}---\textbar{}\allowbreak{}---\textbar{}\allowbreak{}---\textbar{}\allowbreak{}---\textbar{}\allowbreak{}---\textbar{}\allowbreak{}---\textbar{}\allowbreak{}---\textbar{}\allowbreak{}---\textbar{}\allowbreak{}---\textbar{}\allowbreak{}---\textbar{}\allowbreak{}---\textbar{}\allowbreak{}}\\
\texttt{\textbar{}\allowbreak{} sin$^2$2$\theta$ (relic,\allowbreak{} server)\allowbreak{} \textbar{}\allowbreak{} 8.7e-9 \textbar{}\allowbreak{} 2.5e-9 \textbar{}\allowbreak{} 1.2e-9 \textbar{}\allowbreak{} 4.8e-10 \textbar{}\allowbreak{} 2.5e-10 \textbar{}\allowbreak{} 1.4e-10 \textbar{}\allowbreak{} 6.6e-11 \textbar{}\allowbreak{} 3.9e-11 \textbar{}\allowbreak{} 3.3e-11 \textbar{}\allowbreak{} 2.0e-11 \textbar{}\allowbreak{} 1.08e-11 \textbar{}\allowbreak{} 8.8e-12 \textbar{}\allowbreak{} 7.6e-12 \textbar{}\allowbreak{}}\\
\textbf{Points on the relic curve.} The five-point run used sin$^2$2$\theta$ = 2.545e-10, 3.946e-11, 2.022e-11, 1.084e-11 and 8.81e-12 at m\_s = 7.1, 20, 29, 41 and 46 keV. The radiative lifetimes are 1.6e27, 5.82e25, 1.77e25, 5.85e24 and 4.05e24 s. The tool labels three of these five points ''overcloses'' even though Omega\_s h$^2$ = 0.12. This appears to be a strict or rounded inequality in its threshold, not physical overclosure.

\textbf{Fractions at fixed mixing (server).} The fractions f = Omega\_s h$^2$/0.12 are:

\begin{itemize}
  \item At 7.1 keV and sin$^2$2$\theta$ = 5e-11: f = 0.196 (Omega\_s h$^2$ = 0.0236).
  \item At sin$^2$2$\theta$ = 1e-12: f = 0.025 (20 keV), 0.030 (22 keV), 0.049 (29 keV), 0.092 (41 keV) and 0.114 (46 keV).
  \item At sin$^2$2$\theta$ = 1e-11: ten times these values, with 46 and 50 keV overclosing (Omega\_s h$^2$ = 0.136 and 0.158).
\end{itemize}
\textbf{Other outputs.} The free-streaming lengths are 151 kpc at 7.1 keV, 107 kpc at 10 keV and 53.7 kpc at 20 keV. The server provides no momentum distribution or ?p/T?, and no per-point Lyman-alpha flag.

[[FIGURE 10]]

Figure 10 is the server's DW exclusion plot from the 80 $\times$ 80 scan over m\_s = 1--100 keV and sin$^2$2$\theta$ = 1e-14 to 1e-6 (6400 points). Of these, 3092 points overclose and 968 are BBN-excluded. The plot overlays:

\begin{itemize}
  \item the overclosure region,
  \item the Omega\_DM line,
  \item NuSTAR/XMM limits,
  \item a Lyman-alpha thermal-WDM bound,
  \item the 3.5 keV point.
\end{itemize}
Two cautions apply when reading it. First, the provenance of the overlaid X-ray and Lyman-alpha curves is not documented by the tool. Second, the tool's built-in X-ray limits have no coverage at 29--46 keV, so the absence of an X-ray boundary there is a gap in coverage, not an allowed region. The relic line drawn is also the server's normalization, which the next subsection shows sits 10--16x too low in mixing.

\paragraph{Cross-check of the DW yield against the literature}
\textbf{The literature fit.} Abazajian 2005, Eq. 7 (read in the text), gives

m\_s = 3.40 keV (sin$^2$2$\theta$/1e-8)\textasciicircum{}\{-0.615\} (Omega\_DM/0.26)\textasciicircum{}\{0.5\} $\times$ \{0.527 erfc[-1.15 (T\_QCD/170 MeV)\textasciicircum{}\{2.15\}]\}.

The braced factor equals 1 at T\_QCD = 170 MeV, and the fit is valid for 135 \textless{} T\_QCD \textless{} 300 MeV. Inverted, it gives Omega\_s $\propto$ (sin$^2$2$\theta$)\textasciicircum{}\{1.23\} m\_s$^2$. The server instead scales linearly in mixing and roughly as m\textasciicircum{}1.8. The peak production temperature is T\_peak \textasciitilde{} 130 MeV (m\_s/3 keV)\textasciicircum{}\{1/3\}. The critic marked this finding as contested.

\textbf{Relic mixings compared.} The fit gives the relic mixings below (hand arithmetic, Omega\_DM = 0.26, T\_QCD = 170 MeV). Using Omega\_DM = 0.264 changes them by about 1\%.

\begin{longtable}{@{}p{0.240\textwidth}p{0.240\textwidth}p{0.240\textwidth}p{0.240\textwidth}@{}}
\toprule
\textbf{m\_s (keV)} & \textbf{Literature relic sin$^2$2$\theta$} & \textbf{Server relic sin$^2$2$\theta$} & \textbf{Ratio (literature / server)} \\ \midrule \endhead
3 & 1.23e-8 & 1.2e-9 & 10.2 \\
7.1 & 3.02e-9 & 2.5e-10 & 12.1 \\
20 & 5.61e-10 & 3.9e-11 & 14.4 \\
29 & 3.06e-10 & 2.0e-11 & 15.3 \\
41 & 1.75e-10 & 1.08e-11 & 16.1 \\
46 & 1.45e-10 & 8.8e-12 & 16.4 \\
\bottomrule\end{longtable}
So the server overproduces DW sterile neutrinos at fixed mixing. Its relic mixings should be raised by about 10--16 to match this fit. An earlier critic estimate of a 2.5x discrepancy understated the gap.

\textbf{Revised fractions.} Scaling as f = (sin$^2$2$\theta$/sin$^2$2$\theta$\_relic)\textasciicircum{}\{1.23\} (hand arithmetic) gives:

\begin{itemize}
  \item At 7.1 keV: f = 0.0064 for sin$^2$2$\theta$ = 5e-11 and 0.0098 for 7e-11. This is consistent with the literature statement that DW at the 3.5 keV-line mixing makes about 1\% of the dark matter. The server gives 0.196 and 0.27.
  \item At sin$^2$2$\theta$ = 1e-12: f = 4.2e-4 (20 keV), 8.8e-4 (29 keV), 1.75e-3 (41 keV) and 2.2e-3 (46 keV). These are 50--60x below the server values.
  \item At 1e-11: f = 7e-3 (20 keV) to 0.037 (46 keV), so nothing overcloses.
\end{itemize}
These revised fractions are contested. The critic pointed out that at fixed mass the fraction should scale linearly with mixing relative to the literature relic mixing. Applying the 1.23 exponent to the mixing ratio biases f low by a factor of about 2.6--4. The direction of the conclusion is unchanged either way: the server's fractions are far too high.

The fit was also checked only against the 2005 result. Asaka, Laine \& Shaposhnikov (abstract only) find yields that ''differ moderately'' and allow DW as all of the dark matter only below 3.5 keV (6 keV at most). The updated yields in arXiv:1705.01837 and Venumadhav+2016 were not read.

\textbf{Delta N\_eff.} On the relic curve, the sterile number density is fixed by Omega\_s h$^2$/m\_s, whatever yield normalization is used. A hand estimate, assuming one Weyl species, gives Delta N\_eff $\approx$ 11.3 eV/m\_s while the sterile neutrinos are relativistic:

\begin{itemize}
  \item 1.6e-3 at 7.1 keV (the server's on-curve value there is 1.574e-3),
  \item 5.6e-4 at 20 keV,
  \item 2.5e-4 at 46 keV.
\end{itemize}
The earlier server-based statement that Delta N\_eff is ''at most 3e-4'' therefore holds only above about 40 keV. All of these values are 2--3 orders of magnitude below BBN and Planck sensitivity (about 0.2--0.3), so BBN does not constrain keV DW dark matter. A ?p/T? factor of about 0.9 would change the estimate by about 10\%, which does not affect this conclusion.

\paragraph{X-ray constraints at the DW-equivalent masses}
\textbf{Server predictions.} At the server's relic mixings, the predicted line energies are E = 3.55, 10, 14.5, 20.5 and 23 keV for m\_s = 7.1, 20, 29, 41 and 46 keV. The whole-halo M31 fluxes are 1.75e-3, 1.71e-2, 3.88e-2, 8.31e-2 and 0.107 ph/s/cm$^2$.

The tool returns exclusion verdicts but never the limit values themselves:

\begin{itemize}
  \item At 7.1 keV it reports exclusion by NuSTAR M31 (Roach+2020) and XMM blank sky (Dessert+2020).
  \item At 20 keV and sin$^2$2$\theta$ = 3.946e-11 it reports ''EXCLUDED by: NuSTAR M31 (Roach+ 2020)'', with an M31 flux of 0.0171 ph/s/cm$^2$.
  \item At 29 keV and sin$^2$2$\theta$ = 2.022e-11 it reports ''Not excluded by tabulated X-ray limits (within 30\% in mass)'' with no excluding limits listed. That point has $\tau$ = 1.77e25 s and a Milky Way centre whole-halo flux of 264 ph/s/cm$^2$.
\end{itemize}
The 29 keV result is a coverage gap, not an allowed region. The server-based exclusion statements are contested and rest on verdict strings alone.

The critic also flagged the tool's label. Roach+2020 (arXiv:1908.09037) is a Galactic-bulge analysis, while the NuS

\end{document}
